% =====================================================================
%  CONCLUSIONES (versión de revisión, se compila con main_revision.tex)
%
%  Correcciones del 4 de noviembre. Un párrafo por objetivo específico, que es lo
%  que pide la plantilla. Cada uno dice lo obtenido, su causa y lo que implica.
%  No se mencionan defectos ya corregidos ni se repite lo que el documento ya dijo.
% =====================================================================
\chapter*{Conclusiones}
\addcontentsline{toc}{chapter}{CONCLUSIONES}

Al modelar la arquitectura funcional se comprobó que el cambio de piso puede
resolverse como un problema de coordinación, sin que ningún agente suba escaleras ni
use el ascensor. En ninguna misión de las dos campañas falló la asignación ni el
relevo, y a ello contribuyó escribir la planificación sin dependencia de ROS~2, lo que
permitió probar todas las combinaciones de origen y destino sin usar el simulador. Con
un solo agente admisible por piso, la asignación resulta trivial; la dificultad y el aporte del trabajo están en coordinar dos agendas que dependen una de otra durante el relevo.

La plataforma AWS DeepRacer sirve como agente de guiado en interiores, pero su
desempeño depende del sitio y de la tracción. Al no tener codificadores en las ruedas,
su avance solo puede medirse donde el pasillo ofrece estructura al láser: en el
edificio Gregorio la odometría registró como máximo el 73~\% del avance real, mientras
que en Catalayud, con el mismo estimador, cumplió el criterio de la prueba. El sitio de operación debe elegirse como parte del diseño. Por debajo de 0,40~m/s la
tracción no responde, por lo que el agente no tiene una aproximación fina y la
tolerancia de llegada tuvo que fijarse en 0,5~m, el doble de la admitida en simulación.
Con una tarjeta de dos núcleos fue necesario llevar la coordinación a un computador
portátil, y la duración de las sesiones queda limitada por las baterías, que en las pruebas se agotaron en los agentes y en ese computador.

Para operar el sistema basta una página web servida en la red local: en las pruebas
con los agentes físicos cada misión se pidió, se siguió y se confirmó desde un
teléfono, sin instalar nada y sin salida a internet. Sin embargo, las funciones que más
la acercan a una operación real no estaban en el anteproyecto y se añadieron al ver el
sistema en uso. Gracias a la confirmación de cambio de piso, una misión no se da por
completada con el usuario en el piso equivocado, un caso que la instrumentación habría
registrado como éxito, y la cancelación devolvió al agente a su escalera en las seis
misiones canceladas. Por lo tanto, los requisitos de una interfaz como esta no pueden cerrarse de antemano.

Con la evaluación experimental se puede afirmar, con 95~\% de confianza, que la tasa de
éxito del sistema supera el 70~\%, tanto en las treinta misiones en simulación como en
las nueve con los agentes físicos, aunque esas muestras no permiten precisar cuánto la
supera. En ese último caso, la cota se apoya en la localización del propio agente, pues la
llegada no se midió con flexómetro. Como todos los fallos fueron de navegación, sobre
todo de precisión de llegada, y ninguno de asignación ni de relevo, la mejora del
sistema pasa primero por la navegación.
